\documentclass[a4paper,12pt]{report}

\usepackage[utf8]{inputenc}
\usepackage[T1]{fontenc}
\usepackage[french]{babel}
\usepackage{geometry}
\usepackage{graphicx}
\usepackage{hyperref}

\geometry{margin=2.5cm}

\title{CompanionIO, Conversational AI BOT}
\author{Votre Nom}
\date{\today}

\begin{document}

\maketitle

\tableofcontents
\listoffigures
\listoftables

% =====================
% YOUR CONTENT STARTS HERE
% =====================

\chapter*{Introduction Générale}
\addcontentsline{toc}{chapter}{Introduction Générale}

L’intelligence artificielle s’est progressivement imposée comme une composante essentielle des systèmes informatiques modernes. Elle est aujourd’hui intégrée dans une grande variété d’applications, allant des moteurs de recommandation aux systèmes de recherche, en passant par les outils de génération de contenu et les assistants conversationnels.
\setlength{\parskip}{6pt}z

\setlength{\parindent}{0pt}
Cette intégration massive de l’IA dans les applications logicielles a profondément modifié la manière dont les utilisateurs interagissent avec les systèmes numériques. Les fonctionnalités intelligentes ne sont plus des modules isolés, mais des briques intégrées au cœur des architectures applicatives.
\setlength{\parskip}{6pt}

Cependant, malgré ces avancées, l’usage actuel de l’intelligence artificielle reste principalement centré sur des cas d’usage spécifiques et isolés. Les systèmes existants proposent des assistants spécialisés ou des outils orientés tâche, mais ils ne constituent pas encore de véritables assistants numériques autonomes capables de s’intégrer de manière continue dans les activités quotidiennes des utilisateurs.
\setlength{\parskip}{6pt}

En effet, la majorité des solutions actuelles nécessitent une interaction explicite de l’utilisateur et ne disposent pas d’une capacité d’orchestration globale des tâches quotidiennes. Il existe donc un écart entre la puissance des modèles d’IA modernes et leur intégration dans des systèmes capables de réduire significativement la charge cognitive et opérationnelle des utilisateurs.
\setlength{\parskip}{6pt}

Dans ce contexte, ce projet de fin d’études s’inscrit dans une démarche visant à explorer la conception d’un système d’assistance intelligent capable de converser,d’orchestrer et d'executer certaines tâches et d’interagir de manière plus fluide avec l’environnement numérique de l’utilisateur.
% =====================
% END CONTENT
% =====================
\input{chapters/chapter1}

\end{document}